The two pictures on the next page show a small asterism, as seen directly on
the sky and as a mirror image. Three stars are labelled: S1, S2 and Sx. The
position of the asterism is also marked with a rectangle on the larger-scale
map of the sky.

% FIGURE NEEDED: two pictures of the asterism (direct view and mirror image)
% with stars S1, S2, Sx labelled; also marked on the large-scale sky map
% (2011_NIGHT.pdf)

\begin{itemize}
\item Find this asterism and point your telescope to it.
\item Using the illuminated reticle as a fixed reference point, and the
  stopwatch, measure the time taken for the stars S1, S2 and Sx to move across
  the field. You may rotate the eyepiece so that the cross-hairs of the
  reticle are in the most convenient position for your measurement.
\item Use your measurements and the known declinations of stars S1 and S2 as
  given below to determine the declination of star Sx.
\item On the answer sheet, give your measurements and working, and estimate
  the random error in your result.
\end{itemize}

For each set of measurements you make, draw the view through the eyepiece on
the answer sheet. (Use the blank circular field on the answer sheet.)

Mark the drawing with the compass directions N and E. Draw the reticle and the
tracks of the stars to show the motion which you timed using the stopwatch.

Mark the ends of each timed track and show which time measurement refers to
which track --- for example, for measurement ``T1'' marking the ends
``Start T1'' and ``End T1''.

The angle of the reticle can be easily adjusted by rotating the eyepiece
around its optical axis. If you change the angle of the reticle for a new
measurement, draw a new diagram.

The declinations of the field stars S1 and S2 are:
\[ \mathrm{S1}: \ \delta = +19^\circ 48' 18'' \qquad
   \mathrm{S2}: \ \delta = +20^\circ 06' 10'' \]

Assume that: $\delta(\mathrm{S2}) > \delta(\mathrm{Sx}) >
\delta(\mathrm{S1})$.
